\chapter{Topos Theory}
\label{ch:topos}

\section{Introduction}

A topos is a category that behaves like the category $\Set$ of sets, but in a generalized and often constructive sense. Toposes arise naturally as categories of sheaves on topological spaces (or more general sites), as the classifying categories of geometric theories, and as models of constructive set theory.

The theory of toposes was initiated by Grothendieck in the context of algebraic geometry (as ``Grothendieck toposes'') and developed by Lawvere and Tierney into elementary topos theory, which uses no large cardinal axioms and has deep connections to logic and type theory.

\section{Subobject Classifiers}

The most striking feature distinguishing $\Set$ from an arbitrary category is the existence of a \emph{truth-value object} $\Omega$: a set $\Omega = \{T, F\}$ such that every subset $S \subseteq A$ is classified by its characteristic function $\chi_S: A \to \Omega$.

\begin{definition}[Subobject classifier]
In a category $\cat{C}$ with a terminal object $1$, a \textbf{subobject classifier} is an object $\Omega$ together with a morphism $\mathrm{true}: 1 \to \Omega$ such that for every monomorphism $m: S \hookrightarrow A$ in $\cat{C}$, there is a unique \textbf{characteristic morphism} $\chi_m: A \to \Omega$ making the square
\[
  \begin{tikzcd}
    S \arrow[r] \arrow[d, "m"', hook] & 1 \arrow[d, "\mathrm{true}"] \\
    A \arrow[r, "\chi_m"'] & \Omega
  \end{tikzcd}
\]
a pullback. In other words, monomorphisms (subobjects) $m: S \hookrightarrow A$ are in natural bijection with morphisms $A \to \Omega$.
\end{definition}

\begin{example}
In $\Set$: $\Omega = \{0,1\}$ (or $\{F,T\}$), $\mathrm{true} = 1 \mapsto 1$. The characteristic function of $S \subseteq A$ is $\chi_S(a) = 1$ iff $a \in S$.
\end{example}

\begin{example}
In the presheaf category $[\cat{C}\op, \Set]$: $\Omega$ is the presheaf of \textbf{sieves}, $\Omega(c) = \{S \mid S \text{ is a sieve on } c\}$ (a sieve on $c$ is a set of morphisms with codomain $c$ closed under precomposition). This is a generalized truth-value object.
\end{example}

\section{Elementary Toposes}

\begin{definition}[Elementary topos]
An \textbf{elementary topos} (or simply \textbf{topos}) is a category $\mathcal{E}$ that:
\begin{enumerate}
  \item Has all finite limits.
  \item Has all finite colimits. (This follows from the other axioms.)
  \item Has a subobject classifier $\Omega$.
  \item Has \textbf{power objects}: for every object $B \in \mathcal{E}$, there is an object $\Omega^B$ (the ``power object'' of $B$) and a natural bijection
  \[
    \mathcal{E}(A \times B, \Omega) \;\cong\; \mathcal{E}(A, \Omega^B).
  \]
  Equivalently, $\mathcal{E}$ is \textbf{locally cartesian closed}: the functor $- \times B: \mathcal{E} \to \mathcal{E}$ has a right adjoint $(-)^B$ for every $B$.
\end{enumerate}
\end{definition}

\begin{remark}
The power object axiom says $\mathcal{E}$ is \textbf{cartesian closed}: it has exponential objects $B^A = \mathcal{E}(A, B)$ representing the functor $\mathcal{E}(- \times A, B)$.

The combination of finite limits, cartesian closure, and a subobject classifier is remarkably powerful: it suffices to interpret all of higher-order intuitionistic logic.
\end{remark}

\begin{theorem}[Basic consequences]
In any topos $\mathcal{E}$:
\begin{enumerate}
  \item $\mathcal{E}$ has all finite colimits.
  \item $\mathcal{E}$ is locally cartesian closed: every slice category $\mathcal{E}/A$ is itself a topos.
  \item The subobject functor $\mathrm{Sub}: \mathcal{E}\op \to \Set$ is representable by $\Omega$ (via: $\mathrm{Sub}(A) \cong \mathcal{E}(A, \Omega)$).
  \item $\mathrm{Sub}(A)$ carries a Heyting algebra structure (the internal logic is intuitionistic).
\end{enumerate}
\end{theorem}

\section{Examples of Toposes}

\begin{example}[$\Set$]
The category of sets $\Set$ is the archetypal topos. Its subobject classifier is $\Omega = \{0,1\}$ and its internal logic is classical Boolean logic.
\end{example}

\begin{example}[Presheaf toposes]
For any small category $\cat{C}$, the presheaf category $\widehat{\cat{C}} = [\cat{C}\op, \Set]$ is a topos. These are called \textbf{presheaf toposes} (or Grothendieck toposes over a trivial topology). The subobject classifier $\Omega(c)$ is the set of sieves on $c$ in $\cat{C}$; finite limits and colimits are pointwise; exponentials are computed by the formula $F^G(c) = [\cat{C}\op/c, \Set](G|_{c}, F|_{c})$.

Special cases:
\begin{itemize}
  \item $\cat{C} = \mathbf{1}$: $[\mathbf{1}\op, \Set] \cong \Set$.
  \item $\cat{C} = BG$ (a group): $[BG\op, \Set] \cong G$-$\Set$ (category of $G$-sets).
  \item $\cat{C} = \Delta\op$ (simplices): $[\Delta, \Set]$ is the category of simplicial sets.
\end{itemize}
\end{example}

\begin{example}[Sheaf toposes]
Let $X$ be a topological space. The category $\mathbf{Sh}(X)$ of sheaves of sets on $X$ is a topos. The subobject classifier is the sheaf $\Omega$ with $\Omega(U) = \{V \mid V \subseteq U, V \text{ open}\}$ (open subsets). The internal logic of $\mathbf{Sh}(X)$ reflects the topology of $X$: classical logic holds iff $X$ is a discrete space; Heyting (intuitionistic) logic holds in general.
\end{example}

\begin{example}[The effective topos]
The \textbf{effective topos} (Hyland 1982) is a topos whose internal logic is the logic of recursive realizability. It provides a categorical model of Kleene's realizability interpretation and contains a natural number object with the property that all its functions are recursive. This topos cannot be a Grothendieck topos.
\end{example}

\section{The Internal Language of a Topos}

Every topos has an \textbf{internal language}: a formal system (Mitchell-Bénabou language) in which one can reason about objects of the topos using logical language, with the topos playing the role of a model.

\begin{definition}[Mitchell-Bénabou language]
Given a topos $\mathcal{E}$, the \textbf{internal language} $L(\mathcal{E})$ is a typed higher-order language where:
\begin{itemize}
  \item \textbf{Types} correspond to objects of $\mathcal{E}$.
  \item \textbf{Terms of type $A$ in context $\Gamma$} correspond to morphisms $\|\Gamma\| \to A$ in $\mathcal{E}$.
  \item \textbf{Propositions (of type $\Omega$)} correspond to morphisms $\|\Gamma\| \to \Omega$.
  \item \textbf{Logical connectives} ($\wedge, \vee, \Rightarrow, \neg, \forall, \exists$) correspond to morphisms constructed from $\Omega, \times, +, (-)^\Omega$, etc.
  \item \textbf{Subsets} $\{x: A \mid \phi(x)\}$ correspond to pullbacks along $\phi: A \to \Omega$ of $\mathrm{true}: 1 \to \Omega$.
\end{itemize}
\end{definition}

\begin{theorem}[Soundness and completeness]
A formula is provable in the internal language (using intuitionistic logic) if and only if it holds in every topos. The soundness theorem says that valid constructions in the internal language correspond to morphisms in $\mathcal{E}$.
\end{theorem}

\begin{example}
In $\mathbf{Sh}(X)$, the internal statement ``every function $f: \mathbb{R} \to \mathbb{R}$ is continuous'' is satisfied (where $\mathbb{R}$ is the sheaf of continuous real-valued functions). This reflects the fact that the internal real numbers of $\mathbf{Sh}(X)$ are continuous functions.
\end{example}

\section{Heyting Algebras and Intuitionistic Logic}

\begin{definition}[Heyting algebra]
A \textbf{Heyting algebra} is a lattice $H$ with a binary operation $\Rightarrow$ (implication) satisfying $a \leq b \Rightarrow c$ iff $a \wedge b \leq c$ (the residuation property). Equivalently, it is a cartesian closed poset.

A \textbf{Boolean algebra} is a Heyting algebra satisfying $a \vee \neg a = 1$ (law of excluded middle), where $\neg a = a \Rightarrow 0$.
\end{definition}

\begin{theorem}
In any topos $\mathcal{E}$, the subobject lattice $\mathrm{Sub}(A)$ is a Heyting algebra for every $A \in \mathcal{E}$. The topos has Boolean logic (satisfies the law of excluded middle internally) if and only if $\Omega \cong \Omega + \Omega$ (the subobject classifier decomposes as a coproduct of two copies of the terminal object).
\end{theorem}

\begin{corollary}
$\Set$ is a Boolean topos. Sheaf toposes $\mathbf{Sh}(X)$ for non-discrete $X$ are non-Boolean. The effective topos is non-Boolean.
\end{corollary}

\section{Geometric Morphisms}

\begin{definition}[Geometric morphism]
Let $\mathcal{E}$ and $\mathcal{F}$ be toposes. A \textbf{geometric morphism} $f: \mathcal{E} \to \mathcal{F}$ is an adjunction $f^* \dashv f_*$ where:
\begin{itemize}
  \item $f^*: \mathcal{F} \to \mathcal{E}$ is the \textbf{inverse image functor} (left adjoint).
  \item $f_*: \mathcal{E} \to \mathcal{F}$ is the \textbf{direct image functor} (right adjoint).
\end{itemize}
and $f^*$ preserves finite limits (is left exact).
\end{definition}

\begin{remark}
The condition that $f^*$ is left exact (preserves finite limits) is the key additional condition that distinguishes geometric morphisms from arbitrary adjunctions.
\end{remark}

\begin{example}
A continuous map $g: X \to Y$ of topological spaces induces a geometric morphism $g: \mathbf{Sh}(X) \to \mathbf{Sh}(Y)$ via:
\begin{itemize}
  \item $g^*(\mathcal{F})(U) = \colim_{V \supseteq g(U), V \text{ open}} \mathcal{F}(V)$ (sheaf-theoretic inverse image).
  \item $g_*(\mathcal{G})(V) = \mathcal{G}(g^{-1}(V))$ (direct image).
\end{itemize}
\end{example}

\begin{example}[The global sections geometric morphism]
For any topos $\mathcal{E}$, there is a unique geometric morphism $\Gamma: \mathcal{E} \to \Set$ (the \textbf{global sections morphism}):
\begin{itemize}
  \item $\Gamma^*: \Set \to \mathcal{E}$ is the \textbf{constant sheaf functor} (sends a set to the constant presheaf, sheafified if necessary).
  \item $\Gamma_*: \mathcal{E} \to \Set$, $\mathcal{F} \mapsto \mathcal{E}(1, \mathcal{F})$ is the \textbf{global sections functor}.
\end{itemize}
\end{example}

Geometric morphisms are the correct notion of ``morphism'' between Grothendieck toposes, and they form the objects and morphisms of the 2-category of toposes.

\section{Lawvere-Tierney Topologies}

Inside a topos, one can define an internal notion of ``coverage'' or ``topology'' using the subobject classifier.

\begin{definition}[Lawvere-Tierney topology]
A \textbf{Lawvere-Tierney topology} (or \textbf{local operator}) on a topos $\mathcal{E}$ is a morphism $j: \Omega \to \Omega$ satisfying:
\begin{enumerate}
  \item $j \circ \mathrm{true} = \mathrm{true}$ (truth is preserved).
  \item $j \circ j = j$ (idempotence).
  \item $j \circ \wedge = \wedge \circ (j \times j)$ (meets are preserved).
\end{enumerate}
where $\wedge: \Omega \times \Omega \to \Omega$ is the conjunction morphism.
\end{definition}

\begin{theorem}
Given a Lawvere-Tierney topology $j$ on $\mathcal{E}$, the full subcategory of \textbf{$j$-sheaves} (objects $F$ for which $\mathcal{E}(-, F)$ inverts $j$-dense monomorphisms) is itself a topos $\mathbf{Sh}_j(\mathcal{E})$, and the inclusion $\mathbf{Sh}_j(\mathcal{E}) \hookrightarrow \mathcal{E}$ has a left adjoint (the \textbf{$j$-sheafification functor}), making it a geometric morphism.
\end{theorem}

\begin{example}
The double-negation topology $j = \neg\neg: \Omega \to \Omega$ on $\Set$ (or any Boolean topos) is trivial—every presheaf is a sheaf. For presheaf toposes $[\cat{C}\op,\Set]$, the double-negation sheaves form a Boolean topos.
\end{example}

\section{Classifying Toposes}

One of the most profound aspects of topos theory is the correspondence between geometric theories and classifying toposes.

\begin{definition}[Geometric theory]
A \textbf{geometric theory} $\mathbb{T}$ is a theory (in the sense of logic) whose axioms can be expressed using $\top$, $\bot$, $\wedge$ (finite conjunctions), $\vee$ (arbitrary disjunctions), $\exists$ (existential quantification), and $=$ (equality)—but no $\forall$ or $\Rightarrow$.
\end{definition}

\begin{theorem}[Classifying topos]
For any geometric theory $\mathbb{T}$, there exists a topos $\mathbf{Set}[\mathbb{T}]$ (the \textbf{classifying topos} of $\mathbb{T}$) such that for any Grothendieck topos $\mathcal{E}$,
\[
  \{\text{geometric morphisms } \mathcal{E} \to \mathbf{Set}[\mathbb{T}]\} \;\cong\; \{\mathbb{T}\text{-models in } \mathcal{E}\},
\]
naturally in $\mathcal{E}$. The topos $\mathbf{Set}[\mathbb{T}]$ ``classifies'' models of $\mathbb{T}$.
\end{theorem}

\begin{example}
\begin{itemize}
  \item The classifying topos of the theory of objects (no axioms) is $[\mathbf{FinSet}\op, \Set]$ (presheaves on finite sets).
  \item The classifying topos of the theory of local rings is the Zariski topos, relevant to algebraic geometry.
  \item The classifying topos of the theory of linear orders is related to the simplex category $\Delta$.
\end{itemize}
\end{example}

\section{Connections to Logic and Type Theory}

The \textbf{Curry-Howard-Lambek correspondence} unifies:
\begin{center}
\begin{tabular}{ccc}
\toprule
Logic & Type Theory & Category Theory \\
\midrule
Proposition & Type & Object \\
Proof & Term & Morphism \\
Conjunction & Product type & Product \\
Disjunction & Sum type & Coproduct \\
Implication & Function type & Exponential \\
True / False & Unit / Empty & Terminal / Initial \\
Universal & Dependent product $\Pi$ & Right adjoint \\
Existential & Dependent sum $\Sigma$ & Left adjoint \\
\bottomrule
\end{tabular}
\end{center}

\begin{theorem}
The internal language of a topos (equipped with a natural number object) models Martin-Löf dependent type theory (in its impredicative form). Conversely, the category of types and terms of a suitable dependent type theory forms a topos.
\end{theorem}

This correspondence is the foundation of \textbf{homotopy type theory} (HoTT) and \textbf{univalent foundations}, where the category-theoretic perspective (via $(\infty,1)$-toposes) provides semantic models for the type theory.

\section{Exercises}

\begin{exercise}
Show that the subobject classifier $\Omega$ in $[\cat{C}\op, \Set]$ is the functor sending each object $c \in \cat{C}$ to the set of sieves on $c$. Verify the universal property: subobjects of $F$ correspond to natural transformations $F \to \Omega$.
\end{exercise}

\begin{exercise}
Show that in any topos, the subobject lattice $\mathrm{Sub}(A)$ has joins: the join $m_1 \vee m_2$ of two subobjects $m_1: S_1 \hookrightarrow A$ and $m_2: S_2 \hookrightarrow A$ is the image of the induced morphism $S_1 \sqcup S_2 \to A$.
\end{exercise}

\begin{exercise}
Show that every topos $\mathcal{E}$ has a natural number object $N$ (an initial algebra for the functor $1 + -$), and that $N$ supports primitive recursion.
\end{exercise}

\begin{exercise}[$\star$]
Describe the subobject classifier for the topos of $G$-sets $[BG\op, \Set]$ (for a group $G$). What are the ``truth values'' in this topos? (They are the subgroups of $G$.)
\end{exercise}

\begin{exercise}[$\star$]
Prove that a geometric morphism $f: \mathcal{E} \to \mathcal{F}$ with $f^*$ exact (preserving finite limits) induces a morphism of Heyting algebras $\Omega_\mathcal{F} \to \Omega_\mathcal{E}$ (via $f^*(\Omega_\mathcal{F}) \cong \Omega_\mathcal{E}$... actually prove that $f^*(\Omega_\mathcal{F})$ is a subobject classifier in $\mathcal{E}$).
\end{exercise}

\begin{exercise}[$\star\star$]
Let $j: \Omega \to \Omega$ be a Lawvere-Tierney topology on a presheaf topos $[\cat{C}\op, \Set]$. Show that $j$ corresponds to a Grothendieck topology on $\cat{C}$, and that the $j$-sheaves are exactly the sheaves for this Grothendieck topology.
\end{exercise}
